% from ega1-9-fr.tex
\oldpage[I]{178}
lorsque $Y$ est l'image fermée de $X$ par $f$, $f$ est un
\emph{épimorphisme} dans la catégorie des $S$-schémas (T, 1.1). Nous
démontrerons au chap. V que, réciproquement, si l'image fermée $Y'$ de $X$
par $f$ existe et si $f$ est un épimorphisme de $S$-schémas, alors on a
nécessairement $Y'=Y$.
\end{remark}

\begin{proposition}[9.5.8]
\label{I.9.5.8-fr}
Supposons vérifiées les hypothèses de (\hyperref[I.9.5.1-fr]{9.5.1}), et
soit $Y'$ l'image fermée de $X$ par $f$. Pour tout ouvert $V$ de $Y$, soit
$f_V:f^{-1}(V)\to V$ la restriction de $f$~; alors l'image fermée de
$f^{-1}(V)$ par $f_V$ dans $V$ existe et est égale au préschéma induit par
$Y'$ sur l'ouvert $V\cap Y'$ de $Y'$ (autrement dit, au sous-préschéma
$\inf(V,Y')$ de $Y$ (\hyperref[I.4.4.3-fr]{4.4.3})).
\end{proposition}

\begin{proof}
Posons $X'=f^{-1}(V)$~; comme l'image directe de $\mathscr{O}_{X'}$ par
$f_V$ n'est autre que la restriction de $f_*(\mathscr{O}_X)$ à $V$, il est
clair que le noyau $\mathscr{J}'$ de l'homomorphisme
$\mathscr{O}_V\to(f_V)_*(\mathscr{O}_{X'})$ est la restriction de
$\mathscr{J}$ à $V$, d'où aussitôt la proposition.
\end{proof}

On notera que ce résultat s'interprète en disant que la formation de l'image
fermée commute à une extension $Y_1\to Y$ du préschéma de base, qui est une
\emph{immersion ouverte}. Nous verrons au chap. IV qu'il en est de même pour
une extension $Y_1\to Y$ qui est un morphisme \emph{plat}, pourvu que $f$
soit séparé et quasi-compact.

\begin{proposition}[9.5.9]
\label{I.9.5.9-fr}
Soit $f:X\to Y$ un morphisme tel que l'image fermée $Y'$ de $X$ par $f$
existe.
\begin{enumerate}
  \item[(i)] Si $X$ est réduit, il en est de même de $Y'$.
  \item[(ii)] Si on suppose vérifiées les hypothèses de
    (\hyperref[I.9.5.1-fr]{9.5.1}) et si $X$ est irréductible (resp.
    intègre), il en est de même de $Y'$.
\end{enumerate}
\end{proposition}

\begin{proof}
Par hypothèse, le morphisme $f$ se factorise en
$X\xrightarrow{g}Y'\xrightarrow{j}Y$, où $j$ est l'injection canonique.
Comme $X$ est réduit, $g$ se factorise en
$X\xrightarrow{h}Y'_{\mathrm{red}}\xrightarrow{j'}Y'$, où $j'$ est
l'injection canonique (\hyperref[I.5.2.2-fr]{5.2.2}), et il résulte alors de
la définition de $Y'$ que $Y'_{\mathrm{red}}=Y'$. Si de plus les conditions
de (\hyperref[I.9.5.1-fr]{9.5.1}) sont remplies, il résulte de
(\hyperref[I.9.5.4-fr]{9.5.4}) que $f(X)$ est dense dans $Y'$~; si $X$ est
irréductible, il en est donc de même de $Y'$
(\hyperref[0.2.1.5-fr]{0, 2.1.5}). L'assertion relative aux préschémas
intègres résulte de la conjonction des deux autres.
\end{proof}

\begin{proposition}[9.5.10]
\label{I.9.5.10-fr}
Soit $Y$ un sous-préschéma d'un préschéma $X$, tel que l'injection canonique
$i:Y\to X$ soit un morphisme quasi-compact. Il existe alors un plus petit
sous-préschéma fermé $\overline{Y}$ de $X$ majorant $Y$~; son espace
sous-jacent est l'adhérence de celui de $Y$~; ce dernier est ouvert dans son
adhérence, et le préschéma $Y$ est induit sur cet ouvert par $\overline{Y}$.
\end{proposition}

\begin{proof}
Il suffit d'appliquer (\hyperref[I.9.5.1-fr]{9.5.1}) à l'injection $j$, qui
est séparée (\hyperref[I.5.5.1-fr]{5.5.1}) et quasi-compacte par hypothèse~;
(\hyperref[I.9.5.1-fr]{9.5.1}) prouve donc l'existence de $\overline{Y}$ et
(\hyperref[I.9.5.4-fr]{9.5.4}) montre que son espace sous-jacent est
l'adhérence de $Y$ dans $X$~; comme $Y$ est localement fermé dans $X$, il
est ouvert dans $\overline{Y}$, et la dernière assertion provient de
(\hyperref[I.9.5.8-fr]{9.5.8}) appliqué à un ouvert $V$ de $X$ tel que $Y$
soit fermé dans $V$.
\end{proof}

Avec ces notations, si l'injection $V\to X$ est quasi-compacte, et si
$\mathscr{J}$ est le faisceau quasi-cohérent d'idéaux de
$\mathscr{O}_X|V$ définissant le sous-préschéma fermé $Y$ de $V$, il résulte
de (\hyperref[I.9.5.1-fr]{9.5.1}) que le faisceau quasi-cohérent d'idéaux de
$\mathscr{O}_X$ définissant $\overline{Y}$ est le prolongement canonique
(\hyperref[I.9.4.1-fr]{9.4.1}) $\overline{\mathscr{J}}$ de $\mathscr{J}$,
car c'est évidemment le plus grand sous-faisceau d'idéaux quasi-cohérent de
$\mathscr{O}_X$ induisant $\mathscr{J}$ sur $V$.
